\section{Conclusion and Future Work}
\label{sec:conclusion}

We presented \textsc{Mem0-Cognitive}, an \emph{emotion-weighted retention + asynchronous consolidation + adaptive reweighting} framework for LLM long-term memory. The framework is mechanism-level rather than architecture-level: its three modules --- a retention law (\S\ref{subsec:forgetting}, Eq.~\ref{eq:retention}), an asynchronous schema-style consolidation pass (\S\ref{subsec:consolidation}), and a top-$k$-weighted adaptive parameter tuner (\S\ref{subsec:meta_learning}) --- each have a single code-level entry point in the companion repository and can be ablated independently through configuration flags.

Consistent with the artifact-first stance described in the introduction, this version of the paper does \emph{not} report headline numeric results; the empty cells in \S\ref{sec:experiments} are forward references to outputs that an evaluation harness, not shipped with this fork, would produce once implemented. The ablation matrix, metric definitions, and baseline set are fixed in this writeup, so the specific hypotheses a later version would refute or confirm are already part of the public record. We believe this positioning is the honest one while no cognitive evaluation harness has been built.

\paragraph{Limitations.} Several limitations follow directly from the mechanism-level scope we adopted:
\begin{itemize}
    \item \textbf{Emotion extractor is not a validated measurement instrument.} As discussed in \S\ref{subsec:emotion_extractor}, the two-path extractor has not been calibrated against any gold-annotated emotion-intensity corpus, and we report no inter-annotator agreement or cross-model consistency numbers. The retention law's regression properties hold for any extractor satisfying the $E_i \in [0,1]$ contract, but treating the extractor as a psychological measurement would be unsupported.
    \item \textbf{No belief update / generative model.} We deliberately do not position the framework as an instance of active-inference or any other cognitive-architectural theory; we do not implement belief updates, a generative model of user relevance, or an explicit prediction-error signal. The adaptive parameter tuner is a top-$k$-weighted heuristic, not Gaussian-Process Bayesian Optimisation. Replacing it with a proper GP-BO step is a clean drop-in extension and is labelled future work.
    \item \textbf{Clustering algorithm is simple.} The consolidation engine's cluster-identification step is greedy compatible-cluster assignment over the embedding space, not DBSCAN or HDBSCAN. We have not measured the impact of swapping in a stronger clustering routine.
    \item \textbf{Cold-start sensitivity.} The adaptive tuner requires initial interaction data to converge; population-default priors mitigate but do not eliminate this issue.
    \item \textbf{LLM dependency.} Both the emotion extractor's primary path and the consolidation engine's summarisation path rely on LLM quality. The lexicon fallback and the deterministic keep-best fallback, respectively, provide robustness at the cost of extractor / consolidation fidelity.
    \item \textbf{Evaluation scope.} The planned experiments focus on English dialogues; cross-cultural variation in emotional expression and memory patterns is out of scope for this paper.
\end{itemize}

\paragraph{Future Directions.}
\begin{enumerate}
    \item \textbf{Gold-annotated emotion calibration} of the two-path extractor, with reported inter-annotator agreement and cross-model consistency, before the extractor is presented as anything other than a software engineering primitive.
    \item \textbf{Real GP-BO parameter tuner.} Replace the top-$k$-weighted heuristic with a Gaussian-Process surrogate and Expected-Improvement acquisition; the current code exposes the update step as a single private method precisely to make this substitution a minimal change.
    \item \textbf{Stronger clustering.} DBSCAN / HDBSCAN / agglomerative variants of the cluster-identification step, with a controlled ablation against the current greedy compatible-cluster assignment.
    \item \textbf{Multi-modal emotion integration.} Extending the extractor to prosodic (pitch, tempo) and visual features for richer salience modelling in spoken or video-mediated dialogue.
    \item \textbf{Distilled cognitive critic.} Training a lightweight specialised model (e.g.\ 7B parameters) to replace LLM calls for emotion scoring and summarisation, trading off fidelity for cost and latency.
    \item \textbf{Cross-lingual evaluation.} Probing cognitive mechanisms across diverse languages and cultural contexts once the CognitiveBench generator supports localisation.
\end{enumerate}

\paragraph{Broader Impact.} \textsc{Mem0-Cognitive} has potential applications in long-term companionship AI, mental-health support systems, and educational tutors requiring persistent user modelling. The ability to selectively forget also raises ethical considerations: malicious actors could exploit forgetting mechanisms to erase evidence of harmful interactions. We advocate for transparent audit logs (which the consolidation engine already produces) and user-controlled forgetting policies in any production deployment built on this framework.

Code, benchmark configurations, and the reproducibility checklist supporting this paper will be released at an anonymised companion repository linked in the submission materials; the URL is omitted here to preserve double-blind review.
